% ──────────────────────────────────────────────────────────────
\subsection{Distributional conservation law}\label{sec:conservation-add}

We now prove the conservation law entirely within the
distributional framework, avoiding the somewhat formal manipulation
of derivatives of delta functions that appears in
Weinberg~\cite[\S2.8]{Weinberg72}.

\begin{theorem}[Conservation of $T^{\alpha\beta}$]%
\label{thm:conservation}
  If each particle is free ($\dd p_j^\alpha/\dd\tau = 0$), then
  \begin{equation}\label{eq:conservation}
    \partial_\beta\, T^{\alpha\beta} = 0
  \end{equation}
  in the distributional sense.  More generally, in the presence of
  external forces,
  \begin{equation}\label{eq:conservation-force}
    \partial_\beta\, T^{\alpha\beta}(x) = G^\alpha(x)\,,
  \end{equation}
  where the \textbf{force density}
  \[
    G^\alpha(x) = \sum_j \int\!\dd\tau\;
      \frac{\dd p_j^\alpha}{\dd\tau}(\tau)\,
      \delta^{(4)}\!\bigl(x - \gamma_j(\tau)\bigr)
  \]
  is the distributional pushforward of the equations of motion.
\end{theorem}

\begin{proof}
  We prove~\eqref{eq:conservation};
  equation~\eqref{eq:conservation-force} follows by the same
  argument, now retaining the nonzero right-hand side
  of~\eqref{eq:cons-ibp} when $\dd p_j^\alpha/\dd\tau \neq 0$.
  For any $\phi \in \mathcal{D}(\Rn{4})$, the standard definition
  of the distributional derivative (integration by parts with a
  sign) gives:
  \begin{align}
    \langle \partial_\beta T^{\alpha\beta},\,\phi\rangle
      &= -\langle T^{\alpha\beta},\,\partial_\beta\phi\rangle
      \notag\\
      &= -\sum_j \int_{\R} p_j^\alpha(\tau)\,
        \dot\gamma_j^\beta(\tau)\,
        \partial_\beta\phi\bigl(\gamma_j(\tau)\bigr)\;\dd\tau
      \notag\\
      &= -\sum_j \int_{\R} p_j^\alpha(\tau)\,
        \frac{\dd}{\dd\tau}\bigl[\phi(\gamma_j(\tau))\bigr]
        \;\dd\tau\,,
        \label{eq:cons-chain}
  \end{align}
  where the last step uses the chain rule
  $\frac{\dd}{\dd\tau}\phi(\gamma(\tau))
  = \dot\gamma^\beta(\tau)\,\partial_\beta\phi(\gamma(\tau))$.
  Since $p_j^\alpha = m_j\dot\gamma_j^\alpha$ and each~$\gamma_j$
  is smooth (Definition~\ref{def:timelike-worldline}), the
  four-momentum $p_j^\alpha(\tau)$ is smooth, so integration by
  parts is justified.  Since $\gamma_j$ is proper
  (Proposition~\ref{prop:Tmunu-well-defined}), the set
  $\gamma_j^{-1}(\operatorname{spt}\phi)$ is compact, so the
  function $\tau \mapsto \phi(\gamma_j(\tau))$ has compact support
  in~$\R$.  The boundary terms therefore vanish, giving:
  \begin{equation}\label{eq:cons-ibp}
    \langle \partial_\beta T^{\alpha\beta},\,\phi\rangle
      = \sum_j \int_{\R}
        \frac{\dd p_j^\alpha}{\dd\tau}(\tau)\,
        \phi\bigl(\gamma_j(\tau)\bigr)\;\dd\tau\,.
  \end{equation}
  For free particles, $\dd p_j^\alpha/\dd\tau = 0$, and the
  right-hand side vanishes for all~$\phi$.
\end{proof}

\begin{remark}
  The elegance of the distributional proof is worth
  emphasising.  In the $3{+}1$ approach, the conservation law
  requires computing
  $\partial_i\bigl[p^\alpha\,(\dd\gamma^\beta/\dd t)\,
  \delta^{(3)}(\vb{x} - \boldsymbol{\gamma}(t))\bigr]
  + \partial_0\bigl[p^\alpha\,(\dd\gamma^0/\dd t)\,
  \delta^{(3)}(\vb{x} - \boldsymbol{\gamma}(t))\bigr]$,
  which involves the multivariable chain rule
  \[
    \frac{\dd}{\dd t}\,\delta^{(3)}\!\bigl(\vb{x}
    - \boldsymbol{\gamma}(t)\bigr)
    = -\frac{\dd\gamma^i}{\dd t}\,
      \frac{\partial}{\partial x^i}\,
      \delta^{(3)}\!\bigl(\vb{x}
      - \boldsymbol{\gamma}(t)\bigr)
  \]
  to convert time derivatives into spatial derivatives.  The
  distributional proof~\eqref{eq:cons-chain}--\eqref{eq:cons-ibp}
  bypasses this entirely: the chain rule on test functions and
  integration by parts do all the work.
\end{remark}

\begin{exercise}\label{ex:conservation-collision}
  Extend Theorem~\ref{thm:conservation} to particles that undergo
  strictly localised collisions (so that worldlines begin and end
  at discrete spacetime events).  Show that $\partial_\beta
  T^{\alpha\beta} = 0$ provided that four-momentum is conserved
  at each collision vertex: $\sum_{n\in c} p_n^\alpha = 0$ for
  each collision~$c$.
  \emph{Hint:} integrate by parts on each worldline segment;
  the boundary terms at each collision vertex sum to zero by
  momentum conservation.
\end{exercise}
